\documentclass[11pt,a4paper]{article}

\usepackage[margin=2.2cm]{geometry}
\usepackage{amsmath,amssymb}
\usepackage{booktabs}
\usepackage{graphicx}
\usepackage{pdflscape}
\usepackage{enumitem}
\usepackage{xcolor}
\usepackage[hidelinks]{hyperref}
\usepackage{microtype}

\setlist{nosep,leftmargin=1.7em}
\setlength{\parindent}{0pt}
\setlength{\parskip}{0.55em}
\newcommand{\code}[1]{\texttt{\detokenize{#1}}}

\title{Balanced Treecode Parameters Across Mesh Scales}
\author{}
\date{}

\begin{document}
\maketitle

\section{Objective}

The purpose of this experiment is to choose a balanced combination of the
quadrature algebraic accuracy $p$, Cartesian multipole order $P$, and MAC
threshold $\theta$ at each mesh scale. Accuracy is assessed on a common set of
786,432 level-5 target points, using the analytic potential as the reference.
The target is
\[
  \lVert V_{\mathrm{tree}}-V_{\mathrm{ex}}\rVert_{L^2(\Omega)}
  \le 10^{-4}.
\]
All production timings use the CPU backend with 128 OpenMP threads. The
reported values are single wall-clock measurements.

\section{Parameter-selection procedure}

\begin{enumerate}
  \item Screen $p=2,4,6,8$, extending to $p=12,16$ only when needed.
  \item Screen coupled $(P,\theta)$ candidates rather than evaluating the full
        Cartesian product.
  \item Use native leaf-barycenter errors only for preliminary screening.
  \item Validate promising configurations for both adaptive and uniform meshes
        at the same 786,432 target coordinates.
  \item Compute the final error against the analytic solution, not against the
        stored uniform potential values.
  \item Increase parameters selectively if a native-grid winner fails the
        common validation target.
\end{enumerate}

The screened $(P,\theta)$ pairs were
\[
\begin{split}
 &(2,0.9),(2,0.7),(2,0.5),(4,0.7),(4,0.5),(4,0.3),\\
 &(6,0.7),(6,0.5),(6,0.3),(8,0.7),(8,0.5).
\end{split}
\]
The very strict value $\theta=0.1$ was excluded from the staged search because
earlier tests showed a large increase in direct interactions with negligible
accuracy improvement over $\theta=0.3$.

\section{Final balanced configurations}

\begin{landscape}
\begin{table}[p]
\centering
\caption{Balanced parameters, wall-clock times, and common-grid errors.}
\label{tab:balanced-parameters}
\scriptsize
\setlength{\tabcolsep}{3.7pt}
\begin{tabular}{llrrrrrrrrr}
\toprule
Scale & Mesh & $n$ & $p$ & $P$ & $\theta$
& $T_{\mathrm{list}}$ & $T_{\mathrm{moment}}$
& $T_{\mathrm{eval}}$ & $T_{\mathrm{total}}$
& Common-grid $L^2$ error \\
& & & & & & (s) & (s) & (s) & (s) & \\
\midrule
L2 & Adaptive & 1,032   & 16 & 4 & 0.7 & 0.0595 & 0.0141 & 0.0682 & 0.2972 & $4.4614\times10^{-3}$ \\
L2 & Uniform  & 1,536   & 16 & 4 & 0.7 & 0.0225 & 0.0159 & 0.0640 & 0.2968 & $4.4599\times10^{-3}$ \\
L3 & Adaptive & 10,496  & 16 & 8 & 0.7 & 0.0769 & 0.1388 & 0.6308 & 1.3423 & $4.4554\times10^{-4}$ \\
L3 & Uniform  & 12,288  & 16 & 8 & 0.7 & 0.0789 & 0.1509 & 0.6418 & 1.4176 & $4.4559\times10^{-4}$ \\
L4 & Adaptive & 100,432 & 8  & 8 & 0.7 & 1.0107 & 0.3958 & 2.7325 & 6.6935 & $8.3734\times10^{-5}$ \\
L4 & Uniform  & 98,304  & 16 & 8 & 0.7 & 0.8763 & 2.0391 & 10.0355 & 16.9804 & $7.1433\times10^{-5}$ \\
L5 & Adaptive & 823,672 & 8  & 6 & 0.5 & 8.5547 & 3.0858 & 37.1372 & 67.9263 & $9.0704\times10^{-5}$ \\
L5 & Uniform  & 786,432 & 8  & 6 & 0.5 & 7.8554 & 2.4889 & 34.0349 & 65.8349 & $5.9447\times10^{-5}$ \\
\bottomrule
\end{tabular}
\end{table}
\end{landscape}

\section{Main observations}

\begin{itemize}
  \item Neither L2 nor L3 reaches the $10^{-4}$ target under the tested limit
        $p\le16$. They should therefore be described as the best tested coarse
        configurations, not as target-satisfying choices.
  \item Both L4 configurations meet the target. The adaptive mesh needs only
        $p=8$, whereas the uniform mesh needs $p=16$. The adaptive L4 run is
        about 2.54 times faster in total and 3.67 times faster in potential
        evaluation.
  \item Both L5 configurations meet the target with
        $(p,P,\theta)=(8,6,0.5)$. At this scale, the uniform run is about 3.1\%
        faster and has the smaller common-grid error. Adaptivity is therefore
        most advantageous near the intermediate 100,000-element scale in this
        manufactured problem.
  \item The balanced quadrature order decreases as the mesh is refined. Coarse
        tetrahedra require high-order quadrature to resolve the narrow analytic
        Gaussian; fine tetrahedra resolve it geometrically and need a lower
        $p$.
  \item Native leaf-barycenter errors can be misleading. For example, the L3
        adaptive native error is $8.29\times10^{-5}$, but its common-grid error
        is $4.46\times10^{-4}$. Common-target validation is essential for a
        fair adaptive-versus-uniform comparison.
  \item The initial L5 choice $(6,6,0.5)$ also looked adequate on its native
        target set, but failed common validation. Increasing to $p=8$ reduced
        the common error below the target.
\end{itemize}

\section{Reproducibility files}

\begin{itemize}
  \item \code{run_balanced_tuning.sh}: resumable staged-search driver.
  \item \code{balanced_parameter_screen.tsv}: all preliminary screening runs.
  \item \code{balanced_parameter_selection.tsv}: provisional native-grid choices.
  \item \code{balanced_common_validation.tsv}: common-grid validation and corrections.
  \item \code{balanced_final_production.tsv}: final production timing measurements.
  \item \code{balanced_final_results.tsv}: final joined timing and error table.
  \item \code{experiment_results/balanced_parameter_screen/}: screening logs.
  \item \code{experiment_results/balanced_common_validation/}: validation logs.
  \item \code{experiment_results/balanced_final_production/}: final timing logs.
\end{itemize}

For publication-quality timing statistics, repeat each final configuration at
least three times and report the median together with the observed variation.

\end{document}
